\section{Clearing denominators}\label{sec:integerization}

This section proves \eqref{eq:integerizer} for all \(m\in4\mathbb N\). For
\(F\in\mathbb Q[T]\), the Gauss valuation \(v_p(F)\) is the minimum
valuation of its coefficients. Set \(v_p(0)=+\infty\).

\subsection{Denominators and a formal identity in the parameter}

Let \(a\) be a variable for now. Put \(c_i=\binom{i+g}{g}\),
\(U_i=c_i a^i\), \(h_i=\sum_{s=1}^g1/(i+s)\), and
\(Q_s(a)=(a/(a-1))^s-1\). In Section~\ref{sec:source} replace \(2^i\) by
\(a^i\) and \(2^s-1\) by \(Q_s(a)\). This gives the rational coefficient
functions of \(D_gM\) in Lemma~\ref{lem:source-bank}. They follow from
Leibniz's rule and

\[
 \begin{aligned}
 L^{(s)}&=(-1)^{s-1}(s-1)!\bigl((a-1)^{-s}-a^{-s}\bigr),\\
 \beta^{(s)}&=(-1)^{s-1}(s-1)!a^{-s}
          \left(L-\sum_{r=1}^{s-1}\frac{Q_r(a)}r\right).
 \end{aligned}
\]

Put \(U=2m-1\) and \(d_U=\operatorname{lcm}(1,\ldots,U)\). The identity
\begin{equation}\label{eq:int-harmonic}
 c_i h_i=\sum_{s=1}^g\frac{(-1)^{s-1}}s\binom{i+g}{g-s}
\end{equation}
is the derivative at zero of the Vandermonde identity
\(\binom{i+g+x}{g}=\sum_{s=0}^g\binom{i+g}{g-s}\binom{x}{s}\). So the
harmonic denominators need only the indices \(s\le g\le U\). The finite sums
have denominators cleared by \(d_U\) in \(C\) and by \(d_U^2\) in \(B\).
The entries of \(B\) have pole order at most \(g\) at \(a=1\). Scaling rows
gives
\begin{equation}\label{eq:int-Pm}
 P_m(a,T):=d_U^{3m}(a-1)^{gm}G_m(a,T)\in\mathbb Z[a,T].
\end{equation}
Here \(G_m(a,T)\) is the determinant with parameter \(a\). The rest of the
paper uses \(G_m(2,T)\). Its finite rational formulas contain no negative
powers of \(a\).

Work formally in \(\mathbb Q[[a]]\), with \(L_1=\sum_{r\ge1}a^r/r\), \(L_2=\sum_{r\ge1}a^r/r^2\),
and

\[
 \gamma_i=\tfrac12\left(h_i^2-\sum_{s=1}^g(i+s)^{-2}\right).
\]

For all \(g,i\ge0\) the rational functions satisfy
\begin{equation}\label{eq:int-t}
 t_i^{(g)}-L_1U_i=\sum_{k\ne i}\frac{c_ka^k}{i-k}-U_i h_i,
\end{equation}
\begin{equation}\label{eq:int-b}
 b_i^{(g)}+L_1U_i h_i+L_2U_i
       =\sum_{k\ne i}\frac{c_ka^k}{(i-k)^2}+\gamma_iU_i.
\end{equation}
Here and below sums over \(k\) start at zero. To prove these uniformly in
\(g\), adjoin to \(\mathbb Q((a))\) a formal symbol \(\lambda\) with
\(\lambda'=1/a\). Differentiating \(D_g a^x=\binom{x+g}{g}a^x\) in the
auxiliary exponent gives

\[
 D_g(a^i\lambda)=U_i(\lambda+h_i),\quad
 D_g(a^i\lambda^2)=U_i(\lambda^2+2h_i\lambda+2\gamma_i).
\]

The substitutes \(\bar L=-L_1-\lambda\) and
\(\bar\beta=-L_2-\lambda^2/2\) obey the differential equations of \(L\) and
\(\beta\): \(\bar L'=1/(a(a-1))\) and \(\bar\beta'=\bar L/a\). The base
formal identities \(t_i^0-L_1a^i=\sum_{k\ne i}a^k/(i-k)\) and
\(b_i^0+L_2a^i=\sum_{k\ne i}a^k/(i-k)^2\) hold coefficientwise. Apply
\(D_g\); it multiplies each \(a^k\) coefficient by \(c_k\). Then compare the
terms free of the independent symbol \(\lambda\). This proves
\eqref{eq:int-t} and \eqref{eq:int-b}, including the coefficient at
\(k=i\). The formula for \(\gamma_i\) also follows from the coefficient of
\(x^2\) in the same Vandermonde identity.

\subsection{Rank-one features and specialization at 2}

For an independent variable \(X\) form

\[
 \mathcal H(a,X)=\begin{bmatrix}C(a)\\
                B(a)-L_1C(a)+L_2A(a)-XA(a)\end{bmatrix}.
\]

Fix a weight \(k\ge0\). On \(I_m\) put \(e_I(i)=[i=k]\) and
\(w_I(i)=v_I(i)^2\), where \(v_I(i)=1/(i-k)\) for \(i\ne k\) and \(v_I\) is
zero otherwise. Define \(e_J,v_J,w_J\) on \(J_m\) in the same way. By
\eqref{eq:int-t} and \eqref{eq:int-b},
\begin{equation}\label{eq:int-rank-one}
 [a^k]\mathcal H=c_k\left(
 \begin{bmatrix}e_I\\v_I-h_ke_I\end{bmatrix}v_J^t
 +\begin{bmatrix}0\\-e_I\end{bmatrix}w_J^t
 +\begin{bmatrix}v_I-h_ke_I\\-w_I-h_kv_I+(\gamma_k-X)e_I\end{bmatrix}e_J^t
                  \right).
\end{equation}
The second rank-one summand vanishes unless \(k\in I_m\), and the third
unless \(k\in J_m\). So weight \(k\) has at most
\(1+[k\in I_m]+[k\in J_m]\) features. A nonzero determinant term must
choose \(2m\) distinct features. For a fixed coefficient of \(a\) only
finitely many features matter. So finite multilinearity proves this
statement formally.

For \(m=4h\), the cheapest \(8h\) features are two copies of
\(0,\ldots,h-1\) and three of \(h,\ldots,3h-1\). Their total weight is
\begin{equation}\label{eq:int-Nm}
 N_m=2\sum_{k=0}^{h-1}k+3\sum_{k=h}^{3h-1}k
                  =13h^2-4h=\frac{13}{16}m^2-m.
\end{equation}
The determinant-preserving row shear gives
\(\det\mathcal H(a,X)=G_m(a,X-L_2)\). Hence the right side is divisible
by \(a^{N_m}\) in \(\mathbb Q[[a]][X]\). The substitution \(X=T+L_2\) is
polynomial in \(X\), so it preserves divisibility. With the integer
polynomial \eqref{eq:int-Pm} this gives
\begin{equation}\label{eq:int-divisible}
 d_U^{3m}(a-1)^{gm}G_m(a,T)\in a^{N_m}\mathbb Z[a,T].
\end{equation}
Only now set \(a=2\). The factor \((a-1)^{gm}\) is one, so
\begin{equation}\label{eq:int-twoadic}
 v_2(G_m)\ge N_m-3m\lfloor\log_2U\rfloor.
\end{equation}
No value of \(L_1(2)\) or \(L_2(2)\) is used. The formal translation is
undone before the integer specialization. The supplementary reports test
finitely many formal coefficients. These tests only corroborate the
argument, which holds for every \(g\).

\subsection{Valuation bounds at odd primes}\label{subsec:odd}

Let \(p\) be an odd prime with \(p^2>U+g\). Put

\[
 r=g\bmod p,\quad \kappa_n=[n\bmod p\ge p-r],\qquad
 \tau_n=(\kappa_n-1)[n\ge p\ \hbox{or}\ p\le g].
\]

When \(r=0\), \(\kappa_n=0\). Since \(n+g<p^2\), Kummer's carry criterion
gives \(v_p\binom{n+g}{g}=\kappa_n\). If \(p\) divides a denominator index
\(s\), then \(v_p\binom{n+g}{g-s}\ge\kappa_n\); if \(s\le n\), then
\(v_p\binom{n-s+g}{g}=\kappa_n\). All nonzero denominator indices and index
differences have valuation at most one. The ultrametric inequality, applied
to the formulas of Section~\ref{sec:source}, gives the following upper
bounds for negative Gauss valuations.

\begin{center}
\begin{tabular}{p{0.36\textwidth}p{0.54\textwidth}}
\hline
Entry & Upper bound \\
\hline
\(C_{ij}\), \(i\ne j\) & \([i\equiv j\pmod p]-\min(\kappa_i,\kappa_j)\) \\
\(C_{ii}\) & \((1-\kappa_i)[p\le g]\) \\
\(B_{ij}\), \(i\ne j\) & \([i\equiv j\pmod p]-\min(\tau_i,\tau_j)\) \\
\(B_{ii}-TU_i\), \(i\ge p\) or \(2p\le g\) & \(2-\kappa_i\) \\
\(B_{ii}-TU_i\), \(i<p\) and \(p\le g<2p\) & \(1\) \\
\(B_{ii}-TU_i\), \(i<p\) and \(p>g\) & \(0\) \\
\hline
\end{tabular}
\end{center}

For the last three lines, the first diagonal sum has a squared \(p\)
denominator only if \(i\ge p\). Both denominators in its double sum are
divisible by \(p\) only if \(2p\le g\), and then the binomial gains
\(\kappa_i\). A single \(p\) denominator costs at most one. The coefficient
\(-TU_i\) has nonnegative valuation and does not worsen the bounds.

Let \(W_m(p)\) be the maximum weight sum over permutations. By Leibniz,
\(-v_p(G_m)\le W_m(p)\). Integer row and column potentials whose sums
dominate every entry weight also give an exponent, namely their total. No
assignment optimality or determinant rank theorem is needed. The coarse
bound is \(W_m(p)\le3m\).

\subsection{Small primes}

Let \(p<g\). After enlarging the special diagonal bounds if needed, take
\(w_C(i,j)=[i\equiv j\pmod p]-\kappa_i\kappa_j\) and \(w_B=1+w_C\). Let
\(c_I(v),c_J(v)\) be the residue copy counts and \(K_I,K_J\) the killed
index counts. Put
\[
 S_U=\sum_{\kappa_v=0}\min(2c_I(v),c_J(v)),\quad
 S_K=\sum_{\kappa_v=1}\min(2c_I(v),c_J(v)).
\]
Then
\begin{equation}\label{eq:int-W}
 W_m(p)\le\min\{m+S_U,\ 3m-2K_I-K_J+S_K\}.
\end{equation}
For the proof, subtract the bottom-row constant \(m\). The score is then
\(x-z\), where \(x\) counts unkilled same-residue matches and \(z\) killed
off-residue matches. Let \(y\) count killed same-residue matches. Then
\(x\le S_U\) and \(y\le S_K\). Counting the killed rows and columns gives
\(z+y\ge2K_I+K_J-2m+x\). Both bounds follow.

For \(\rho=p/m\) and \(t=v/m\in[0,\rho)\), the copy counts are
\[
 c_I(t)=\#\{n:1/4\le t+n\rho<5/4\},\quad
 c_J(t)=\#\{n:0\le t+n\rho<2\}.
\]
They change only at the six endpoints
\begin{equation}\label{eq:int-endpoints}
 0,\rho,\quad a-\lfloor a/\rho\rfloor\rho\ (a=1/4,5/4,2),
 \quad-1/2+(\lfloor1/(2\rho)\rfloor+1)\rho.
\end{equation}
The last is the killed threshold. Every endpoint is \(b+n\rho\) with
integer \(n\) and \(b\in\{0,1/4,5/4,2,-1/2\}\). So all walls lie among the
numbers \(|b-c|/n\), \(n\ge1\). Each fixed positive strip contains finitely
many of them. Within a wall cell the objectives in \eqref{eq:int-W} are
affine. We integrate their lower envelope after also splitting at the
affine crossings. Each endpoint in \eqref{eq:int-endpoints}, multiplied by
\(m\), is an integer. With half-open conventions the residue counts are
therefore exactly \(m\) times the cell widths, with the correct boundary
assignments. No asymptotic equality to the true assignment profile is used.

Let \(\delta=1/128\). The 756 affine pieces in the certificate specify the
exact integral \(L_*\) on \([\delta,1/2]\). Its value is the following
numerator divided by the denominator.
\begin{center}
\begin{minipage}{0.92\linewidth}\footnotesize\ttfamily
\seqsplit{1448276314740194964855260100800462127260298031042796749529725627370010396734895633658285333350450905704443029571002838710237}
\par\medskip /\par\medskip
\seqsplit{1404098881580803157259376918521625735839461675553873436297217652662216162782957043593966060046854826150978176673015236224000}
\end{minipage}
\end{center}

Near zero the row potentials \(1-\kappa_i\) and \(2-\kappa_i\), with zero
column potentials, are feasible. They give \(W_m(p)\le3m-2K_I\). Any
interval of \(m\) integers contains at least
\(\lfloor m/p\rfloor(g\bmod p)\) killed indices. Hence
\begin{equation}\label{eq:int-near-zero}
 W_m(p)/m\le3-2\{1/(2\rho)\}+2\rho.
\end{equation}
Write \(\delta=1/(2Q)\), \(Q=64\). The substitution \(x=1/(2\rho)\) gives
\[
 \int_0^\delta(3-2\{1/(2\rho)\})d\rho
 =\frac1{2Q}+\sum_{q\ge Q}\left(\frac1q-\log(1+1/q)\right).
\]
The summand exceeds \(1/(2q(q+1))\) by exactly
\(\int_0^1s(1-s)/(q(q+s)(q+1))\,ds\le1/(6q^2(q+1))\). Summing, and adding
the integral of \(2\rho\), gives the rigorous upper bound
\begin{equation}\label{eq:int-small}
 I_{\rm small}\le2\delta+2\delta^2=129/8192.
\end{equation}

\subsection{Large primes and the integerizer}

For \(\rho>1/2\), keep each row copy in both blocks and every column copy.
In a residue cell with flag \(\kappa\), copy digit \(t\) has
\(\tau=(\kappa-1)[t\ge1]\). The exact data give an integer potential to
each row and column type. Every pair of types must satisfy the off-residue
bound of Section~\ref{subsec:odd}. A pair in the same residue cell must also
satisfy the same-residue bound, including the diagonal case. Checking both
cases also covers two distinct residues in the same cell. The objective is
the sum of potentials times the exact cell counts.

The exact rational verifier in the supplementary materials reconstructs
the complete wall geometry. It checks all 756 low pieces and all 1,942 high
integer inequalities. It gives the following high upper profile.

\begin{center}
\begin{tabular}{p{0.36\textwidth}p{0.54\textwidth}}
\hline
Interval in \(\rho\) & Upper profile \\
\hline
\((1/2,7/12)\) & \(1\) \\
\((7/12,5/8)\) & \(3\rho -3/4\) \\
\((5/8,2/3)\) & \(\rho +1/2\) \\
\((2/3,3/4)\) & \(4\rho -3/2\) \\
\((3/4,5/6)\) & \(3\rho -3/4\) \\
\((5/6,1)\) & \(7/4\) \\
\((1,7/4)\) & \(15/4-2\rho\) \\
\((7/4,2)\) & \(2-\rho\) \\
\((2,5/2)\) & \(0\) \\
\hline
\end{tabular}
\end{center}

Its integral is \(287/192\). The copy geometry splits this table into 11
cells. The data list the individual types and integer potentials. No
optimizer or numerical rounding is used. On each wall in a positive strip
use \(3m\) instead. This costs \(O(m\log m)\) in total and covers
\(\rho=1/2\).

Choose the exponents \(E_p\) in the following order.
\begin{enumerate}
\item At 2 take \(E_2=3m\lfloor\log_2U\rfloor-N_m\).
\item At odd \(p^2\le U+g\) use \(3m\lfloor\log_p U\rfloor\).
\item At \(p>U\) use zero.
\item At the remaining \(p/m<\delta\) use \(3m-2K_I\).
\item At a wall in a positive strip use \(3m\).
\item Elsewhere below \(m/2\) use \eqref{eq:int-W}.
\item Elsewhere above \(m/2\) use the least feasible high dual objective.
\end{enumerate}
The positive rational product
\begin{equation}\label{eq:int-Dm}
 D_m=\prod_p p^{E_p}
\end{equation}
is finite, and every exponent bounds \(-v_p(G_m)\). So \(D_m\) integerizes
every coefficient. A negative \(E_2\) only divides out a proved common power
of two.

The positive 2-adic overhead is \(O(m\log m)\). Small odd primes cost at
most \(3m\,\pi(\sqrt{U+g})\log U=O(m^{3/2}\log m)=o(m^2)\). On each fixed
positive strip we use the classical weighted prime number theorem
\[
 \frac1m\sum_p f(p/m)\log p\longrightarrow\int f(\rho)d\rho
\]
for a bounded finitely piecewise smooth function \(f\). For
\eqref{eq:int-near-zero}, first omit \(p<\varepsilon m\). Then apply the
theorem on \([\varepsilon,\delta]\), and let \(\varepsilon\downarrow0\). By
\(3m\theta(\varepsilon m)\), the omitted part has normalized limsup at most
\(3\varepsilon\). The exact counts on positive strips and the finite wall
cost now prove

\[
 \limsup\frac{\log D_m}{m^2}
 \le L_*+\frac{129}{8192}+\frac{287}{192}-\frac{13}{16}\log2.
\]

Finally let \(\ell_-=2\sum_{k=0}^{29}(1/3)^{2k+1}/(2k+1)<\log2\). The same
verifier checks the exact rational comparison

\[
 \begin{aligned}
 \tau_{\rm basic}:={}&L_*+\frac{129}{8192}+\frac{287}{192}
                            -\frac{13}{16}\ell_-\\
 ={}&1.97881984510269917464907037626379\ldots<1.979.
 \end{aligned}
\]

The decimal is the exact rational of the receipt, rounded. This proves
\eqref{eq:integerizer}, and in particular the weaker rate \(\tau\le1.99\).
